\documentclass[conference]{IEEEtran}
\usepackage{amsmath}
\usepackage{booktabs}
\usepackage{graphicx}

\title{Proekt Kladenets: A Dual-Use Synthetic Research Platform for Autonomous Swarm Coordination}

\author{\IEEEauthorblockN{George David Tsitlauri}
\IEEEauthorblockA{\textit{Dept. of Informatics \& Telecommunications}\\
\textit{University of Thessaly, Greece}\\
 gdtsitlauri@gmail.com}}

\begin{document}
\maketitle

\begin{abstract}
Proekt Kladenets is a synthetic proof-of-concept repository for decentralized autonomous-swarm research. It combines perception, communications, swarm-role logic, tracking, explainability, and dashboard-facing outputs inside a single workspace. The current repository contains real module-level artifacts and selected constrained-hardware experiments, but it does not provide real flight, hardware-in-the-loop, or operational validation. The contribution of the repository is therefore architectural breadth and simulation-backed integration rather than a mature field-ready UAV platform.
\end{abstract}

\section{Introduction}
Autonomous swarm systems require tight integration across perception, communication, navigation, decision logic, and human oversight. Proekt Kladenets explores that integration through a synthetic research platform oriented toward low-cost, decentralized operation.

\section{Repository Scope}
The repository includes:
\begin{itemize}
\item module-level perception and spiking-style experiments,
\item communication simulations including cognitive-radio-style behavior and line-of-sight scenarios,
\item swarm-role and deception simulations,
\item navigation, spoofing, and explainability outputs,
\item committed artifact files under \texttt{results/}.
\end{itemize}

\section{Current Evidence}
The committed artifacts show that the repository is more than a conceptual scaffold. Example files include:
\begin{itemize}
\item \texttt{results/ilya\_muromets\_bench.txt} with a small local benchmark,
\item \texttt{results/dobrynya\_rewards.txt} with Q-learning reward traces,
\item \texttt{results/vasilisa\_xai\_iff.txt} with explanation and override output,
\item multiple trajectory, spoofing, and line-of-sight plots.
\end{itemize}
These are useful research artifacts, but they remain simulation-heavy and subsystem-specific.

\section{Interpretation Boundary}
The repository should not be presented as demonstrating operational swarm superiority or field deployment. Its value lies in broad cyber-physical integration, low-cost experimentation, and a transparent synthetic evidence trail.

\section{Future Work}
The natural next steps are hardware-in-the-loop testing, richer communication realism, and narrower task-focused evaluation that goes deeper on fewer subsystems.

\section{Why the Current Release Still Stands Up Well}
Even with its simulation-only scope, the current Kladenets release still stands
up as a meaningful research repository because it integrates many swarm-system
subproblems in one place and already produces committed artifacts across those
modules. Its strength is architectural synthesis and synthetic experimentation,
not field validation, and that is a legitimate contribution when stated
explicitly.

\section{Conclusion}
Proekt Kladenets is a credible synthetic research platform with broad technical scope. It is strongest as an integration-heavy proof of concept and should be presented that way.

\bibliographystyle{IEEEtran}
\bibliography{references}
\end{document}



